\begin{definition}[Restoring bond multiplicities]%
    \label{def:peps_multiplicity_bond_map}
    \lean{TNLean.PEPS.multiplicityBondMap,
        TNLean.PEPS.multiplicityBondMap_apply,
        TNLean.PEPS.multiplicityBondMap_single,
        TNLean.PEPS.multiplicityBondFamilyMap,
        TNLean.PEPS.multiplicityBondFamilyMap_apply}
    \leanok
    Let $V_i$ and $M_i$, indexed by a finite set, be finite-dimensional
    coordinate spaces, with
    $m_i=\dim M_i>0$. On the direct sum of matrix spaces, define
    \begin{align}
        J((X_i)_i) & =\bigl(m_i^{-1/2}X_i\otimes I_{M_i}\bigr)_i.
    \end{align}
    On each matrix unit this reads
    \begin{align}
        \ket{k}\bra{\ell}\longmapsto
        \frac{1}{\sqrt{m_i}}\sum_{a=0}^{m_i-1}
        \ket{k,a}\bra{\ell,a}.
    \end{align}
    For an irreducible representation of dimension $d_i$, the regular
    representation has multiplicity $m_i=d_i$. Thus $J$ is precisely the
    bond transformation used in the semi-regular bond-dimension reduction
    of \cite[Section~7]{Schuch2010PEPS}.
    In a fixed block $i$, the four open indices give
    $J(X_i)_{(k,a),(\ell,b)}=m_i^{-1/2}(X_i)_{k,\ell}\delta_{a,b}$:
    % Formula: J(X_i)_(k,a),(ell,b)=m_i^(-1/2) (X_i)_(k,ell) delta_(a,b).
    % Source: Papers/1001.3807/paper_v3.tex:2982--3013, the two bond
    % transformations figs6/prove-equiv-semireg-reg-left and -right.
    % Ink-to-index: the upper west/east legs are k,ell in V_i; the lower
    % west/east legs are a,b in M_i. All four are virtual bond indices.
    % There is no contraction on either side. The right panel is the tensor
    % product X_i tensor I_{M_i}, so both signatures are (V,P)=(4,0).
    % The block label i is fixed, not an additional open leg or a summed index.
    \begin{tenkzequation}
        \tenkzkernel
        \begin{tenkz}[rows={wire,wire}, cols=1, bonds=none]
            \tn[name=multiplicityImage, skin=box, label pos=180,
                ports={135:virtual:$k$, 45:virtual:$\ell$,
                        225:virtual:$a$, 315:virtual:$b$}]{J(X_i)}
        \end{tenkz}
        $=m_i^{-1/2}$
        \begin{tenkz}[rows={wire,wire}, cols=1, bonds=none]
            \tn[at=(1,1), skin=box,
                ports={180:virtual:$k$, 0:virtual:$\ell$}]{X_i}
            \tn[at=(2,1), skin=box,
                ports={180:virtual:$a$, 0:virtual:$b$}]{I_{M_i}}
        \end{tenkz}
    \end{tenkzequation}
\end{definition}

\begin{theorem}[The multiplicity-restoring bond isometry]%
    \label{thm:peps_multiplicity_bond_isometry}
    \lean{TNLean.PEPS.trace_multiplicityBondMap_conjTranspose_mul,
        TNLean.PEPS.multiplicityBondMap_injective,
        TNLean.PEPS.sum_trace_multiplicityBondFamilyMap_conjTranspose_mul,
        TNLean.PEPS.multiplicityBondFamilyMap_injective}
    \leanok
    \uses{def:peps_multiplicity_bond_map}
    The map $J$ preserves the Hilbert--Schmidt inner product:
    \begin{align}
        \sum_i\tr[J(X)_i^\dagger J(Y)_i]
        =\sum_i\tr[X_i^\dagger Y_i].
    \end{align}
    In particular, it is injective.
\end{theorem}

\begin{proof}\leanok
    For each block,
    \begin{align}
        \tr\bigl[(m_i^{-1/2}X_i\otimes I)^\dagger
            (m_i^{-1/2}Y_i\otimes I)\bigr]
         & =m_i^{-1}\tr[X_i^\dagger Y_i]\tr[I_{M_i}] \\
         & =\tr[X_i^\dagger Y_i].
    \end{align}
    Summing gives the assertion. The Hilbert--Schmidt norm vanishes
    only on the zero matrix, which proves injectivity.
\end{proof}

\begin{theorem}[The weighted semi-regular bond becomes a regular bond]%
    \label{thm:peps_multiplicity_bond_weight}
    \lean{TNLean.PEPS.multiplicityBondMap_sqrt_smul,
        TNLean.PEPS.multiplicityBondFamilyMap_sqrt_smul}
    \leanok
    \uses{def:peps_multiplicity_bond_map}
    For any matrices $D_i$, one has
    \begin{align}
        J\bigl((\sqrt{m_i}D_i)_i\bigr)
        =\bigl(D_i\otimes I_{M_i}\bigr)_i.
    \end{align}
    Taking $m_i=d_i$ and $D_i=D^i(gh^{-1})$ gives the two bond vectors
    compared in \cite[Section~7]{Schuch2010PEPS}. This establishes the
    isometry on the bond matrix spaces; the equality of complete PEPS
    contractions after its application on all bonds is a further statement.
\end{theorem}

\begin{proof}\leanok
    In each block the factors $m_i^{-1/2}$ and $\sqrt{m_i}$ cancel.
\end{proof}
